% ===== 第四部分：Fourier 分析 / 第3章：Fourier级数的一些应用 =====
% !TEX root = ../../main.tex
\setchapterimage{figure/P4/P4C3封面图.png}
\setchapterintro{本章通过三个方向展示 Fourier 级数的应用：利用 Parseval 恒等式证明等周不等式，利用三角多项式与 Fourier 系数建立 Weyl 等分布定理，以及构造连续但处处不可微的函数。这些例子说明 Fourier 方法不仅能够处理函数展开，也能把几何、数论和函数正则性问题转化为对频率成分的研究。}
\chapter{Fourier级数的一些应用}

% 第1节：等周不等式
% !TEX root = ../../../main.tex
% 所属：第四部分：Fourier 分析 / 第3章：Fourier级数的一些应用
% 节：3.1 等周不等式
% 正文按原文件顺序保存；由 ../chapter3.tex 调用。
\section{等周不等式}

这一节中,我们考虑的问题是：对一条有固定周长的闭曲线,其围成面积的最大值.在此之前我们需要一些基本的定义（即便其中的一些在 Part2 中已经讲过了）.

\begin{definition}
	一条参数化的曲线 $\gamma$ 为一个映射
\[
\gamma:[a,b]\to\R^2.
\]
$\gamma$ 的图像为平面上的一个点集,称为一条曲线,记为 $\Gamma$.

若 $\Gamma$ 不是自相交的,称其为简单曲线.

若它的两个端点重合,称其为闭曲线.

在上述的参数化形式下,两种情况可表示为 $\gamma(s_1)\neq\gamma(s_2)$（除非 $s_1=a$ 和 $s_2=b$,此时有 $\gamma(a)=\gamma(b)$）.

\end{definition}
可以将 $\gamma$ 延拓为 $\R$ 上周期为 $b-a$ 的周期函数,因此可以把 $\gamma$ 视为圆上的函数.

对于任意 $C^1$ 的双射 $s:[c,d]\to[a,b]$,通过公式
\[
\eta(t)=\gamma(s(t))
\]
可以给出 $\Gamma$ 的另一种参数式.显然,对于给定的参数式,$\Gamma$ 是否为闭的或是简单的是独立于参数表达式的.称两个参数式 $\gamma$ 和 $\eta$ 是等价的,若对于所有的 $t$,$s'(t)>0$；这意味着 $\gamma$ 和 $\eta$ 在曲线 $\Gamma$ 上诱导了相同的方向.如果 $s'(t)<0$,则 $\eta$ 是反向的.

\begin{definition}
	若 $\Gamma$ 通过 $\gamma(s)=(x(s),y(s))$ 参数化定义,则曲线 $\Gamma$ 的长度定义为
\[
\ell=\int_a^b\abs{\gamma'(s)}\dd s
=\int_a^b\bigl(x'(s)^2+y'(s)^2\bigr)^{1/2}\dd s.
\]
\end{definition}
曲线 $\Gamma$ 的长度是它的固有属性,不依赖于参数式.为表明这一点,假定 $\eta(t)=\gamma(s(t))$.通过变量替换公式与链式法则,有
\[
\int_a^b\abs{\gamma'(s)}\dd s
=\int_c^d\abs{\gamma'(s(t))}\abs{s'(t)}\dd t
=\int_c^d\abs{\eta'(t)}\dd t,
\]
在下面定理的证明中,我们将使用 $\Gamma$ 参数式的一种特殊形式.
\begin{definition}
	称 $\gamma$ 是弧长化的参数式,若对于所有的 $s$,$\abs{\gamma'(s)}=1$.
\end{definition}这意味着,$\gamma(s)$ 以一个常速度变化,从而,$\Gamma$ 的长度精确地为 $b-a$.因此,经过一个可行的变换,一个弧长化的参数式可以在 $[0,\ell]$ 上定义.

\begin{proposition}
	任意曲线都有一个弧长化的参数式.
\end{proposition}

\begin{proof}
	设 $\eta:[a,b]\to\R^2$ 是曲线 $\Gamma$ 的一个 $C^1$ 参数式,且 $\eta'(t)\neq0$.定义弧长函数 $h:[a,b]\to[0,\ell]$
	\[
	h(s)=\int_a^s\abs{\eta'(t)}\dd t.
	\]
	由 $\eta'(t)\neq0$,有 $\abs{\eta'(s)}>0$,从而
	\[
	h'(s)=\abs{\eta'(s)}>0.
	\]
	故 $h(s)$ 单调,存在反函数 $h^{-1}:[0,\ell]\to[a,b]$.定义新的参数化
	\[
	\gamma(u)=\eta\bigl(h^{-1}(u)\bigr),
	\qquad 0\leq u\leq\ell.
	\]
	令 $s=h^{-1}(u)$,则
	\[
	\gamma'(u)=\eta'(s)\bigl(h^{-1}(u)\bigr)'
	=\frac{\eta'(s)}{\abs{\eta'(s)}},
	\]
	从而
	\[
	\abs{\gamma'(u)}=\frac{\abs{\eta'(s)}}{\abs{\eta'(s)}}=1,
	\]
	且 $\gamma$ 与 $\eta$ 有相同的方向.得证.
\end{proof}

现在回到等周问题上来.给出一条简单闭曲线围成的区域 $\mathcal A$ 的面积的准确公式的工作涉及许多艰深的问题.在许多简单情况下,使用Green公式，面积可由如下的积分式
\begin{equation}\label{eq:area-formula}
	\begin{aligned}
		\mathcal A
		&=\frac12\abs{\int_\Gamma(x\,\dd y-y\,\dd x)}\\
		&=\frac12\abs{\int_a^b\bigl(x(s)y'(s)-y(s)x'(s)\bigr)\dd s}.
	\end{aligned}
\end{equation}

\begin{theorem}
	\textbf{（等周不等式）}
	假定 $\Gamma$ 是 $\R^2$ 中一条长为 $\ell$ 的简单闭曲线,记 $\mathcal A$ 为这条曲线围成的区域的面积.则
	\[
	\mathcal A\leq\frac{\ell^2}{4\pi},
	\]
	等号当且仅当 $\Gamma$ 为圆时成立.
\end{theorem}

\begin{proof}
考虑映射 $(x,y)\mapsto(\delta x,\delta y)$,它把平面 $\R^2$ 映射到自身.观察曲线长度的公式,可知若曲线 $\Gamma$ 的长度为 $\ell$,则它在这个映射下的象长度为 $\delta\ell$.因此这一操作将长度扩大或缩小了 $\delta$ 倍（取决于 $\delta>1$ 或 $\delta<1$）.类似地,可以看出这一映射将面积扩大或缩小了 $\delta^2$.

令 $\delta=2\pi/\ell$,则只需证明若 $\ell=2\pi$,则 $\mathcal A\leq\pi$,等号当且仅当 $\Gamma$ 为圆时成立.
	
	令 $\gamma:[0,2\pi]\to\R^2$ 满足 $\gamma(s)=(x(s),y(s))$ 为曲线 $\Gamma$ 的弧长化的参数式,也就是说,$x'(s)^2+y'(s)^2=1$,对所有的 $s\in[0,2\pi]$.这表明
	\begin{equation}\label{eq:isoperimetric-arc}
		\frac1{2\pi}\int_0^{2\pi}\bigl(x'(s)^2+y'(s)^2\bigr)\dd s=1.
	\end{equation}
	因为曲线是闭的,函数 $x(s)$ 和 $y(s)$ 是以 $2\pi$ 为周期的,所以可以考虑其 Fourier 级数
	\[
	x'(s)\sim\sum a_ne^{ins}
	\quad\text{和}\quad
	y'(s)\sim\sum b_ne^{ins}.
	\]
	对式 \eqref{eq:isoperimetric-arc} 应用 Parseval 等式,得到
	\begin{equation}\label{eq:isoperimetric-parseval}
		\sum_{n=-\infty}^{\infty}\abs{n}^2\bigl(\abs{a_n}^2+\abs{b_n}^2\bigr)=1.
	\end{equation}
	现在对 $\mathcal A$ 的定义式应用 \ref{Parseval 等式的双线性形式}.因为 $x(s)$ 和 $y(s)$ 是实值的,故有 $a_n=\overline{a_{-n}}$ 和 $b_n=\overline{b_{-n}}$,因此可知
	\[
	\mathcal A
	=\frac12\abs{\int_0^{2\pi}\bigl(x(s)y'(s)-y(s)x'(s)\bigr)\dd s}
	=\pi\abs{\sum_{n=-\infty}^{\infty}n\bigl(a_n\overline{b_n}-b_n\overline{a_n}\bigr)}.
	\]
	接下来观察到
	\begin{equation}\label{eq:isoperimetric-key-estimate}
		\abs{a_n\overline{b_n}-b_n\overline{a_n}}
		\leq2\abs{a_n}\abs{b_n}
		\leq\abs{a_n}^2+\abs{b_n}^2.
	\end{equation}
	因为 $\abs{n}\leq\abs{n}^2$,用式 \eqref{eq:isoperimetric-parseval} 得
	\[
	\mathcal A
	\leq\pi\sum_{n=-\infty}^{\infty}\abs{n}^2\bigl(\abs{a_n}^2+\abs{b_n}^2\bigr)
	\leq\pi.
	\]
	
	当 $\mathcal A=\pi$ 时,因为 $\abs{n}\geq2$ 时有 $\abs{n}<\abs{n}^2$,从上述讨论中可知
	\[
	x(s)=a_{-1}e^{-is}+a_0+a_1e^{is}
	\quad\text{和}\quad
	y(s)=b_{-1}e^{-is}+b_0+b_1e^{is}.
	\]
	可知 $x(s)$ 和 $y(s)$ 是实值的,所以 $a_{-1}=\overline{a_1}$ 和 $b_{-1}=\overline{b_1}$.式 \eqref{eq:isoperimetric-parseval} 表明 $2(\abs{a_1}^2+\abs{b_1}^2)=1$.式 \eqref{eq:isoperimetric-key-estimate} 中等号成立,则 $\abs{a_1}=\abs{b_1}=1/2$.记
	\[
	a_1=\frac12e^{i\alpha},
	\qquad
	b_1=\frac12e^{i\beta},
	\]
	$1=2\abs{a_1\overline{b_1}-\overline{a_1}b_1}$ 表明 $\abs{\sin(\alpha-\beta)}=1$,因此当 $k$ 为奇数时,$\alpha-\beta=k\pi/2$.从中可得
	\[
	x(s)=a_0+\cos(\alpha+s)
	\quad\text{和}\quad
	y(s)=b_0\pm\sin(\alpha+s),
	\]
	其中 $y(s)$ 表达式中的符号取决于 $(k-1)/2$ 的奇偶性.可以看出无论哪种情况,等号成立仅当 $\Gamma$ 是一个圆.证明完毕.
\end{proof}
\newpage
%
% 第2节：Weyl 等分布定理
% !TEX root = ../../../main.tex
% 所属：第四部分：Fourier 分析 / 第3章：Fourier级数的一些应用
% 节：3.2 Weyl 等分布定理
% 正文按原文件顺序保存；由 ../chapter3.tex 调用。
\section{Weyl 等分布定理}

若 $x$ 为一个实数,记 $[x]$ 为不大于 $x$ 的最大整数,称 $[x]$ 为 $x$ 的整数部分,$x$ 的分数部分则定义为 $\bkt{x}=x-[x]$.

现在考虑一个不为零的实数 $\gamma\neq0$,观察序列 $\gamma,2\gamma,3\gamma,\ldots$. 考虑分数部分的序列
\[
\bkt{\gamma},\bkt{2\gamma},\bkt{3\gamma},\ldots
\]

\begin{proposition}
\begin{enumerate}
\item 如果 $\gamma$ 是有理数,那么只有有限多个不同的分数出现在 $\bkt{n\gamma}$ 中.
\item 如果 $\gamma$ 是无理数,那么 $\bkt{n\gamma}$ 中所有的数都不同.
\end{enumerate}
\end{proposition}

\begin{proof}
\begin{enumerate}
	\item 若 $\gamma=p/q$,则序列的前 $q$ 项为
	\[
	\bkt{p/q},\bkt{2p/q},\ldots,\bkt{(q-1)p/q},\bkt{qp/q}=0,
	\]
	之后序列开始循环前 $q$ 项,因为
	\[
	\bkt{(q+1)p/q}=\bkt{1+p/q}=\bkt{p/q}.
	\]
	\item 假设存在两个相同的数,立即有对于某些 $n_1\neq n_2$,存在 $\bkt{n_1\gamma}=\bkt{n_2\gamma}$. 则 $n_1\gamma-n_2\gamma\in\Z$,因此 $\gamma$ 是有理数.
\end{enumerate}	
\end{proof}

事实上,Kronecker 证明了,如果 $\gamma$ 是无理数,那么 $\bkt{n\gamma}$ 在 $[0,1)$ 中是稠密的. 换句话说,$\bkt{n\gamma}$ 在 $[0,1)$ 的每个区间中都可以取到值.这是本节核心定理的推论.

\begin{definition}
称 $[0,1)$ 中的序列 $\xi_1,\xi_2,\ldots,\xi_n,\ldots$ 是\textbf{等分布}的,若对于任意区间 $(a,b)\subset[0,1)$,
\[
\lim_{N\to\infty}\frac{\#\{1\leq n\leq N:\xi_n\in(a,b)\}}{N}=b-a,
\]
此处 $\#A$ 意为有限集 $A$ 的基数. 这意味着,对于足够大的 $N$,当 $n\leq N$ 时,$\xi_n\in(a,b)$ 的比例等于区间 $(a,b)$ 的长度与区间 $[0,1)$ 的长度之比. 换句话说,序列 $\xi_n$ 平均地扫过了整个区间,每个子集中有相同的分布. 需要注意的是，序列的排序十分重要.
\end{definition}

\begin{theorem}
\textbf{(Weyl 等分布定理)}
若 $\gamma$ 是无理数,则分数部分的序列 $\bkt{\gamma},\bkt{2\gamma},\bkt{3\gamma},\ldots$ 在 $[0,1)$ 中是等分布的.
\label{Weyl等分布定理}
\end{theorem}

特别地,$\bkt{n\gamma}$ 在 $[0,1)$ 中是稠密的. Kronecker 的定理可以作为它的推论.

固定 $(a,b)\subset[0,1)$,令 $\chi_{(a,b)}(x)$ 为 $(a,b)$ 的特征函数,即函数在 $(a,b)$ 为 1,在 $[0,1)-(a,b)$ 为 0. 将其以 1 为周期延拓到 $\R$ 上,记这一延拓仍然记为 $\chi_{(a,b)}(x)$. 容易得到
\[
\#\{1\leq n\leq N:\bkt{n\gamma}\in(a,b)\}=\sum_{n=1}^N\chi_{(a,b)}(n\gamma).
\]
于是定理成立,等价于当 $N\to\infty$ 时,
\[
\frac1N\sum_{n=1}^N\chi_{(a,b)}(n\gamma)\to\int_0^1\chi_{(a,b)}(x)\dd x.
\]

问题的核心在于以下的结果.

\begin{lemma}
若 $f$ 是连续函数,且以 1 为周期,$\gamma$ 为无理数,则当 $N\to\infty$ 时,
\[
\frac1N\sum_{n=1}^N f(n\gamma)\to\int_0^1 f(x)\dd x.
\]
\label{Weyl等分布的核心引理}
\end{lemma}

\begin{proof}
引理的证明分为三步.

\textbf{第一步:}首先验证 $f$ 为指数函数 $1,e^{2\pi ix},\ldots,e^{2\pi ikx},\ldots$ 时极限的存在性. 若 $f=1$,则极限显然存在. 若 $f=e^{2\pi ikx}$ ($k\neq0$),则积分为零. 因为 $\gamma$ 是无理数,有 $e^{2\pi ik\gamma}\neq1$. 因此,当 $N\to\infty$ 时,
\[
\frac1N\sum_{n=1}^N f(n\gamma)=\frac1N\cdot\frac{e^{2\pi ik\gamma}(1-e^{2\pi ikN\gamma})}{1-e^{2\pi ik\gamma}}
\]
趋于 0.

\textbf{第二步:}若 $f$ 和 $g$ 满足引理,显然对于任意 $A,B\in\C$,$Af+Bg$ 也满足引理. 因此,第一步意味着引理对于所有的三角多项式成立.

\textbf{第三步:}令 $\eps>0$. $f$ 为任意周期为 1 的连续函数. 选定三角多项式 $P$ 满足
\[
\sup_{x\in\R}\abs{f(x)-P(x)}<\eps/3
\]
（由 Stone--Weierstrass 定理可知）. 则由第一步,对于足够大的 $N$,有
\[
\abs{\frac1N\sum_{n=1}^N P(n\gamma)-\int_0^1 P(x)\dd x}<\eps/3.
\]
因此
\[
\begin{aligned}
\abs{\frac1N\sum_{n=1}^N f(n\gamma)-\int_0^1 f(x)\dd x}
\leq{}&\frac1N\sum_{n=1}^N\abs{f(n\gamma)-P(n\gamma)}
+\abs{\frac1N\sum_{n=1}^N P(n\gamma)-\int_0^1 P(x)\dd x}\\&+\int_0^1\abs{P(x)-f(x)}\dd x
<\eps;
\end{aligned}
\]
从而引理得证.
\end{proof}
\begin{proof}[定理的证明]
	对任给 $\eps>0$,取阶梯函数 $f^-,f^+$ 使得
	\[
	\chi_{(a+\eps,b-\eps)}\leq f^-\leq\chi_{(a,b)}\leq f^+\leq\chi_{(a-\eps,b+\eps)},
	\]
	\begin{center}
		\includegraphics[width=0.7\linewidth]{"figure/P4/appoarch to chi(x)"}
	\end{center}
	
	从而
	\[
	b-a-2\eps\leq\int_0^1 f^-(x)\dd x
	\quad\text{和}\quad
	\int_0^1 f^+(x)\dd x\leq b-a+2\eps.
	\]
	若 $S_N=\dfrac1N\sum\limits_{n=1}^N\chi_{(a,b)}(n\gamma)$,则有
	\[
	\frac1N\sum_{n=1}^N f^-(n\gamma)\leq S_N\leq\frac1N\sum_{n=1}^N f^+(n\gamma),
	\]
	因此
	\[
	b-a-2\eps\leq\liminf_{N\to\infty}S_N
	\quad\text{和}\quad
	\limsup_{N\to\infty}S_N\leq b-a+2\eps.
	\]
	因为对所有 $\eps>0$ 这都是对的,所以 $\lim\limits_{N\to\infty}S_N$ 存在且等于 $b-a$. 这样就完成了等分布定理的证明.
\end{proof}


\begin{proposition}
引理 \ref{Weyl等分布的核心引理} 的结论对于所有周期为 1 的,在 $[0,1]$ 上 Riemann 可积的函数都成立.
\label{Weyl等分布的Riemann可积推论}
\end{proposition}

\begin{proof}
假设 $f$ 是实值的,考虑 $[0,1]$ 区间的一个划分 $0=x_0<x_1<\ldots<x_N=1$. 定义
\[
f_U(x)=\sup_{x_{j-1}\leq y\leq x_j}f(y),\quad x\in[x_{j-1},x_j),
\]
\[
f_L(x)=\inf_{x_{j-1}\leq y\leq x_j}f(y),\quad x\in(x_{j-1},x_j].
\]
若 $x\in(x_{j-1},x_j)$,则显然有 $f_L\leq f\leq f_U$,且
\[
\int_0^1 f_L(x)\dd x\leq\int_0^1 f(x)\dd x\leq\int_0^1 f_U(x)\dd x.
\]
进一步地,选择合适的划分,可以保证对于给定的 $\eps>0$,
\[
\int_0^1 f_U(x)\dd x-\int_0^1 f_L(x)\dd x\leq\eps.
\]
根据定理, $f_L$ 与 $f_U$ 都是区间的特征函数的有限线性组合,所以
\[
\frac1N\sum_{n=1}^N f_L(n\gamma)\to\int_0^1 f_L(x)\dd x,
\quad
\frac1N\sum_{n=1}^N f_U(n\gamma)\to\int_0^1 f_U(x)\dd x.
\]
由夹逼即得
\[
\frac1N\sum_{n=1}^N f(n\gamma)\to\int_0^1 f(x)\dd x.
\]
\end{proof}

\begin{theorem}
\textbf{(Weyl 准则)}
一个 $[0,1)$ 中的实数序列是等分布的,当且仅当对于所有的整数 $k\neq0$,当 $N\to\infty$ 时,
\[
\frac1N\sum_{n=1}^N e^{2\pi ik\xi_n}\to0.
\]
\end{theorem}

\begin{proof}
\textbf{充分性:} 由 $\dfrac1N\sum\limits_{n=1}^N e^{2\pi ik\xi_n}\to0$,重复引理 \ref{Weyl等分布的核心引理} 的第一步,后面的三角多项式逼近与夹逼证明同理.

\textbf{必要性:} 设 $\xi_n\subset[0,1)$ 等分布,记
\[
A_N(f)=\frac1N\sum_{n=1}^N f(\xi_n).
\]
先证对任意 $f\in C^0$ (周期为 1),都有
\[
A_N(f)\to\int_0^1 f(x)\dd x.
\]

任取 $I=(a,b)\subset[0,1)$. 由 $\xi_n$ 的等分布性,
\[
A_N(\chi_{(a,b)})=\frac{\#\{1\leq n\leq N:\xi_n\in(a,b)\}}{N}\to b-a=\int_0^1\chi_{(a,b)}(x)\dd x.
\]
对单点集 $C=\{a\}$,同样可推出 $\dfrac{\#\{1\leq n\leq N:\xi_n\in C\}}{N}\to0$,从而
\[
A_N(\chi_{\{a\}})\to0=\int_0^1\chi_{\{a\}}(x)\dd x.
\]
因此
\[
A_N(\chi_{[a,b)})=A_N(\chi_{(a,b)})+A_N(\chi_{\{a\}})\to b-a,
\]
端点不影响极限.

下面考虑分段函数
\[
g(x)=\sum_{j=1}^m C_j\chi_{[x_{j-1},x_j)}(x),\qquad 0=x_0<x_1<\ldots<x_m=1.
\]
由线性性,
\[
A_N(g)=\sum_{j=1}^m C_jA_N(\chi_{[x_{j-1},x_j)})\to\sum_{j=1}^m C_j(x_j-x_{j-1})=\int_0^1 g(x)\dd x.
\]

最后,任取 $f\in C^0$ (周期为 1). $f$ 在 $[0,1]$ 上一致连续,故对任意 $\eps>0$,可取一个充分细的分割 $0=x_0<x_1<\ldots<x_m=1$,使得对任意 $x,y\in[x_{j-1},x_j]$ 有 $\abs{f(x)-f(y)}<\eps$. 定义阶梯函数
\[
g(x)=f(x_{j-1}),\qquad x\in[x_{j-1},x_j).
\]
则 $\norm{f-g}_\infty<\eps$. 估计
\[
\abs{A_N(f)-\int_0^1 f\dd x}
\leq\abs{A_N(f-g)}+\abs{A_N(g)-\int_0^1 g\dd x}+\abs{\int_0^1(g-f)\dd x}.
\]
因此
\[
\varlimsup_{N\to\infty}\abs{A_N(f)-\int_0^1 f\dd x}\leq2\eps.
\]
由 $\eps$ 的任意性得
\[
\frac1N\sum_{n=1}^N f(\xi_n)\to\int_0^1 f(x)\dd x.
\]

取 $f(x)=e^{2\pi ikx}$($k\in\Z,k\neq0$),则 $f\in C^\infty$,从而
\[
\frac1N\sum_{n=1}^N e^{2\pi ik\xi_n}\to\int_0^1 e^{2\pi ikx}\dd x=0.
\]
必要性得证.
\end{proof}
%
% 第3节：构造一个连续但是处处不可微的函数
% !TEX root = ../../../main.tex
% 所属：第四部分：Fourier 分析 / 第3章：Fourier级数的一些应用
% 节：3.3 构造一个连续但是处处不可微的函数
% 正文按原文件顺序保存；由 ../chapter3.tex 调用。
\section{构造一个连续但是处处不可微的函数}

这里,我们对 Weierstrass 的一般构造取一个特例,并且用另外方式证明它.

 我们需要先了解一些相关的知识:部分和,Cesàro 和,推延平均.

\begin{definition}
定义\textbf{推延平均}
\[
\Delta_N(g)(x)=2\sigma_{2N}(g)(x)-\sigma_N(g)(x).
\]
\end{definition}

由 Fejér 核 $F_N$ 的卷积表示
\[
\sigma_N(g)=g*F_N,
\]
立即可得
\[ \Delta_N(g)=g*[2F_{2N}-F_N]. \]


若 $g(x)\sim\sum a_ne^{inx}$,则利用上述定义直接计算可得
\begin{enumerate}
\item $S_N$ 在 $|n|\leq N$ 时,对 $a_ne^{inx}$ 贡献 1 倍,在 $|n|>N$ 时,贡献 0 倍;
\item $\sigma_N$ 在 $|n|\leq N$ 时,对 $a_ne^{inx}$ 贡献 $1-\dfrac{|n|}{N}$ 倍,在 $|n|>N$ 时,贡献 0 倍;
\item $\Delta_N$ 在 $|n|\leq N$ 时,对 $a_ne^{inx}$ 贡献 1 倍,在 $N<|n|\leq 2N$ 时,贡献 $2(1-\dfrac{|n|}{2N})$ 倍,在 $|n|>2N$ 时,贡献 0 倍.
\end{enumerate}
\begin{center}
	\includegraphics[width=0.7\linewidth]{"figure/P4/fig for 3.3"}
\end{center}
\begin{theorem}

若 $0<\alpha<1$,则函数
\[
f(x)=f_\alpha(x)=\sum_{n=0}^{\infty}2^{-n\alpha}e^{i2^n x}
\]
连续但是处处不可微.
\end{theorem}

由级数的绝对收敛性,连续性是显然的. $f$ 的关键性质是它的大多数 Fourier 系数为零,像上式那样缺少很多项的 Fourier 级数称为一个\textbf{缺项 Fourier 级数}.

 对于像 $f$ 这样的缺项级数,推延平均与部分和是基本相等的. 特别地,对函数 $f=f_\alpha$ 有

\[ S_N(f)=\Delta_{N'}(f), \]

其中 $N'$ 是不大于 $N$ 的形如 $2^k$ 的最大整数. 观察 $f$ 的定义,这一点是显然的.

下面使用反证法证明,即假设对于某些 $x_0$,$f'(x_0)$ 存在. 先建立一个关于 Cesàro 平均导数的估计.

\begin{lemma}
设 $g$ 为任意在 $x_0$ 处可微的连续函数,则其 Cesàro 平均满足
\[
\sigma_N(g)'(x_0)=O(\log N).
\]
因此 $\Delta_N(g)'(x_0)=O(\log N)$.
\label{Cesàro平均导数的对数增长}
\end{lemma}

\begin{proof}
由卷积公式 $\sigma_N(g)=g*F_N$ 与 $g$ 在 $x_0$ 处的可微性,有
\[
\sigma_N(g)'(x_0)
=\frac1{2\pi}\int_{-\pi}^{\pi}F_N'(t)g(x_0-t)\dd t.
\]
又因为 $F_N$ 的周期性给出 $\displaystyle\int_{-\pi}^{\pi}F_N'(t)\dd t=0$,故可减去常数 $g(x_0)$:
\[
\sigma_N(g)'(x_0)
=\frac1{2\pi}\int_{-\pi}^{\pi}F_N'(t)\bigl[g(x_0-t)-g(x_0)\bigr]\dd t.
\]

假设 $g$ 在 $x_0$ 处可微,则对 $t$ 充分小有 $\abs{g(x_0-t)-g(x_0)}\leq C\abs{t}$. 因此
\[
\abs{\sigma_N(g)'(x_0)}
\leq\frac C{2\pi}\int_{-\pi}^{\pi}\abs{F_N'(t)}\abs{t}\dd t.
\]

现在对 $\abs{F_N'(t)}$ 建立两个估计.
\begin{itemize}
\item 由 Fejér 核的 Fourier 形式 $F_N(t)=\sum\limits_{|k|\leq N}(1-\dfrac{\abs{k}}{N})e^{ikt}$ 求导得
\[
\abs{F_N'(t)}\leq\sum_{|k|\leq N}\abs{k}(1-\dfrac{\abs{k}}{N})\leq AN^2.
\]
\item 利用 $\sin(\dfrac{Nt}{2})\leq CN\abs{t}$ 与 $\sin(\dfrac t2)\geq c\abs{t}$ (对 $\abs{t}\leq\pi$),可以建立
\[
\abs{F_N'(t)}\leq\frac A{\abs{t}^2}.
\]
\end{itemize}

综合两个估计并对积分区间拆分:
\[
\begin{aligned}
\abs{\sigma_N(g)'(x_0)}
&\leq\frac C{2\pi}\left[\int_{\abs{t}\geq1/N}\abs{F_N'(t)}\abs{t}\dd t
+\int_{\abs{t}\leq1/N}\abs{F_N'(t)}\abs{t}\dd t\right]\\
&\leq\frac{CA}{2\pi}\left[\int_{\abs{t}\geq1/N}\frac{\dd t}{\abs{t}}
+N^2\int_{\abs{t}\leq1/N}\abs{t}\dd t\right]\\
&=\frac{CA}{2\pi}\bigl[O(\log N)+O(1)\bigr]=O(\log N).
\end{aligned}
\]

由 $\Delta_N(g)=2\sigma_{2N}(g)-\sigma_N(g)$ 立即得 $\Delta_N(g)'(x_0)=O(\log N)$.
\end{proof}

\begin{lemma}
若 $2N=2^k$ (即 $N=2^{k-1}$),则
\[
\Delta_{2N}(f)(x)-\Delta_N(f)(x)=2^{-k\alpha}e^{i2^k x}.
\]
\label{缺项Fourier级数的差分}
\end{lemma}

\begin{proof}
考虑$f$的定义 ，这一点显然
\end{proof}

\begin{proof}[定理的证明]
连续性由绝对收敛性即得. 只需证 $f$ 在圆周上处处不可微.

设存在 $x_0$,使 $f'(x_0)$ 存在. 则 $f$ 在 $x_0$ 处可微,且 $f$ 显然是连续函数. 由引理 \ref{Cesàro平均导数的对数增长},
\[
\abs{[\Delta_{2N}(f)-\Delta_N(f)]'(x_0)}=O(\log N).
\]

另一方面,任取 $k\geq1$,令 $N=2^{k-1}$,从而 $2N=2^k$. 由引理 \ref{缺项Fourier级数的差分},
\[
\Delta_{2N}(f)(x)-\Delta_N(f)(x)=2^{-k\alpha}e^{i2^k x}.
\]
对 $x$ 求导得
\[
[\Delta_{2N}(f)-\Delta_N(f)]'(x)=i2^k\cdot2^{-k\alpha}e^{i2^k x}.
\]
特别地,
\[
\abs{[\Delta_{2N}(f)-\Delta_N(f)]'(x_0)}=2^k\cdot2^{-k\alpha}=2^{k(1-\alpha)}=2N^{1-\alpha}.
\]

由于 $0<\alpha<1$,故 $1-\alpha>0$,从而 $2N^{1-\alpha}\to\infty$,其增长速度比 $\log N$ 快得多. 这与引理 \ref{Cesàro平均导数的对数增长} 矛盾. 因此 $f'(x_0)$ 不存在,即 $f$ 处处不可微.
\end{proof}
%
